\documentclass{article}
\usepackage{amsmath,amssymb,amsthm}
\usepackage{algorithm}
\usepackage{algpseudocode}
\usepackage{hyperref}
\usepackage{fullpage}
\newtheorem{theorem}{Theorem}
\newtheorem{lemma}{Lemma}
\newtheorem{assumption}{Assumption}
\newcommand{\E}{\mathbb{E}}
\newcommand{\norm}[1]{\left\|#1\right\|}
\newcommand{\inner}[2]{\left\langle #1,#2\right\rangle}
\begin{document}

\section*{Algorithm Name}
\textbf{AdaGrad for Non‑Convex Smooth Optimization (Scalar Version)}  

The method keeps a cumulative sum of past (scalar) gradients and uses its
square‑root as a denominator for the learning rate.

\section*{Mathematical Formulation}
Let $f:\mathbb{R}\rightarrow\mathbb{R}$ be differentiable.
Denote the stochastic gradient at iteration $k$ by $g_k:=\nabla f(w_k;\xi_k)$,
where $\xi_k$ is a random data sample (or mini‑batch).  

Initialize $w_0\in\mathbb{R}$ and set the accumulated squared gradient
\[
G_0 \;:=\; 0 .
\]

For $k=0,1,\dots,T-1$ perform
\begin{align}
G_{k+1} &:= G_k + g_k^{\,2}, \tag{1}\\[4pt]
\eta_{k} &:= \frac{\alpha}{\sqrt{G_{k+1}+\epsilon}}\,, \tag{2}\\[4pt]
w_{k+1} &:= w_k - \eta_{k}\, g_k , \tag{3}
\end{align}
where $\alpha>0$ is a base stepsize and $\epsilon>0$ is a small
stability constant (often $10^{-8}$).  
Equation~(2) yields the adaptive learning rate
$\eta_k = \alpha / \sqrt{\sum_{i=0}^{k} g_i^{2}+\epsilon}$.

\section*{Pseudocode}
\begin{algorithm}[H]
\caption{AdaGrad (scalar)}\label{alg:adagrad}
\begin{algorithmic}[1]
\Require Initial point $w_0$, base stepsize $\alpha>0$, stability $\epsilon>0$
\State $G_0 \gets 0$
\For{$k=0$ \textbf{to} $T-1$}
    \State Sample $\xi_k$ and compute stochastic gradient $g_k \gets \nabla f(w_k;\xi_k)$
    \State $G_{k+1} \gets G_k + g_k^{\,2}$ \Comment{cumulative squared gradient}
    \State $\eta_k \gets \alpha / \sqrt{G_{k+1}+\epsilon}$ \Comment{adaptive stepsize}
    \State $w_{k+1} \gets w_k - \eta_k\, g_k$ \Comment{parameter update}
\EndFor
\State \Return $w_T$
\end{algorithmic}
\end{algorithm}

\section*{Dry Run Example}
Consider the simple quadratic
\[
f(w) = (w-3)^2 .
\]
Its (exact) gradient is $\nabla f(w)=2(w-3)$.  
We use the deterministic version of AdaGrad (i.e. $g_k=\nabla f(w_k)$),
base stepsize $\alpha=0.1$, and $\epsilon=10^{-8}$.
Initialize $w_0=0$ and $G_0=0$.

\begin{center}
\begin{tabular}{c|c|c|c|c}
$k$ & $w_k$ & $g_k=2(w_k-3)$ & $G_{k+1}=G_k+g_k^2$ & $w_{k+1}=w_k-\eta_k g_k$\\ \hline
0 & $0$ & $-6$ & $0+36=36$ & $\eta_0=0.1/\sqrt{36}=0.1/6=0.0166667$,\\
  &   &   &   & $w_1=0-0.0166667(-6)=0.1000$\\ \hline
1 & $0.1$ & $2(0.1-3)=-5.8$ & $36+33.64=69.64$ & $\eta_1=0.1/\sqrt{69.64}=0.1/8.347=0.01197$,\\
  &   &   &   & $w_2=0.1-0.01197(-5.8)=0.1695$\\ \hline
2 & $0.1695$ & $2(0.1695-3)=-5.661$ & $69.64+32.05=101.69$ & $\eta_2=0.1/\sqrt{101.69}=0.1/10.084=0.00992$,\\
  &   &   &   & $w_3=0.1695-0.00992(-5.661)=0.2222$
\end{tabular}
\end{center}
Objective values:
\[
f(w_0)=9,\; f(w_1)= (0.1-3)^2=8.41,\; f(w_2)= (0.1695-3)^2\approx8.28,\;
f(w_3)\approx8.22 .
\]
The adaptive stepsizes shrink as the accumulated squared gradient grows,
illustrating the AdaGrad behaviour.

\section*{Assumptions}
\begin{assumption}[Smoothness]\label{ass:smooth}
$f:\mathbb{R}\rightarrow\mathbb{R}$ is $L$‑smooth, i.e.,
\[
\forall x,y\in\mathbb{R}:\quad
f(y)\le f(x)+\inner{\nabla f(x)}{y-x}+\frac{L}{2}(y-x)^2 .
\]
\end{assumption}

\begin{assumption}[Unbiased Stochastic Gradient]\label{ass:unbiased}
For each $k$, the stochastic gradient satisfies
\[
\E_{\xi_k}[g_k \mid w_k]=\nabla f(w_k).
\]
\end{assumption}

\begin{assumption}[Bounded Variance]\label{ass:variance}
There exists $\sigma^2\ge0$ such that
\[
\E_{\xi_k}\!\big[\|g_k-\nabla f(w_k)\|^{2}\mid w_k\big]\le\sigma^{2},
\qquad\forall k .
\]
\end{assumption}

\begin{assumption}[Gradient Lower Bounded]\label{ass:gradbound}
The stochastic gradients are almost surely bounded:
\[
\|g_k\|\le G_{\max}\quad\text{for all }k .
\]
\end{assumption}

All constants $\alpha,\epsilon,G_{\max},L,\sigma$ are positive and fixed.

\section*{Main Convergence Theorem}
\begin{theorem}[Non‑Convex AdaGrad Convergence]\label{thm:main}
Let $\{w_k\}_{k=0}^{T}$ be generated by Algorithm~\ref{alg:adagrad}
under Assumptions~\ref{ass:smooth}–\ref{ass:gradbound}.
Define $w^{*}$ as a (global) minimizer of $f$ and let
$\Delta_f:=f(w_0)-f(w^{*})\ge0$.
Then the average expected squared gradient satisfies
\[
\frac{1}{T}\sum_{k=0}^{T-1}\E\!\big[\|\nabla f(w_k)\|^{2}\big]
\;\le\;
\frac{2\Delta_f}{\alpha\,\sqrt{T\,\epsilon}}
\;+\;
\frac{2L\,G_{\max}^{2}\,\alpha}{\sqrt{T\,\epsilon}} .
\tag{4}
\]
Consequently,
\[
\min_{0\le k<T}\E\!\big[\|\nabla f(w_k)\|^{2}\big]
= O\!\left(\frac{1}{\sqrt{T}}\right).
\]
\end{theorem}

\section*{Theoretical Properties}
\begin{itemize}
\item \textbf{Stability.} The denominator $\sqrt{G_{k+1}+\epsilon}$ never
vanishes; the algorithm is thus numerically stable even when gradients are
small.
\item \textbf{Hyper‑parameter Sensitivity.}
  The bound~\eqref{eq:4} is decreasing in $\alpha$ up to the
  $2L G_{\max}^{2}\alpha$ term, showing a trade‑off:
  larger $\alpha$ speeds up early progress (first term) but inflates the second term.
\item \textbf{Comparison to SGD.}
  For constant stepsize $\eta$, SGD yields an $O(1/\sqrt{T})$ bound on the
  average gradient norm with a factor proportional to $\eta$. AdaGrad automatically
  reduces the stepsize when the cumulative gradient grows,
  guaranteeing the same asymptotic rate without manual tuning.
\end{itemize}

\section*{Discussion}
The theorem mirrors classical AdaGrad results for convex objectives,
but here we exploit smoothness to control the descent of the function value.
The $1/\sqrt{T}$ decay is optimal for first‑order methods on non‑convex
$L$‑smooth functions when only unbiased stochastic gradients are available.

\section*{Preliminaries (Definitions and Lemmas)}
\begin{lemma}[Smoothness Descent Lemma]\label{lem:smooth}
If $f$ is $L$‑smooth, then for any $x$ and any step $d$,
\[
f(x+d)\le f(x)+\inner{\nabla f(x)}{d}+\frac{L}{2}d^{2}.
\]
\end{lemma}
\begin{proof}
Directly from Definition~\ref{ass:smooth} with $y=x+d$.
\end{proof}

\begin{lemma}[Bound on Cumulative Adaptive Steps]\label{lem:stepsum}
For the scalar AdaGrad stepsizes $\eta_k$ defined in \eqref{2},
\[
\sum_{k=0}^{T-1}\eta_k \;=\;
\alpha\sum_{k=0}^{T-1}\frac{1}{\sqrt{G_{k+1}+\epsilon}}
\;\le\;
\frac{2\alpha}{\sqrt{\epsilon}}\sqrt{T\;G_{T} } .
\tag{5}
\]
\end{lemma}
\begin{proof}
Since $G_{k+1}=G_k+g_k^{2}$, the sequence $\{G_k\}$ is non‑decreasing.
Apply Cauchy–Schwarz:
\[
\Bigl(\sum_{k=0}^{T-1}\frac{1}{\sqrt{G_{k+1}+\epsilon}}\Bigr)^{2}
\le
\Bigl(\sum_{k=0}^{T-1}1\Bigr)\Bigl(\sum_{k=0}^{T-1}\frac{1}{G_{k+1}+\epsilon}\Bigr)
= T\sum_{k=0}^{T-1}\frac{1}{G_{k+1}+\epsilon}.
\]
Because $x\mapsto 1/(x+\epsilon)$ is decreasing,
\[
\sum_{k=0}^{T-1}\frac{1}{G_{k+1}+\epsilon}
\le
\int_{0}^{G_T}\frac{1}{u+\epsilon}\,du
= \log\!\Bigl(1+\frac{G_T}{\epsilon}\Bigr)
\le \frac{G_T}{\epsilon}.
\]
Thus
\[
\sum_{k=0}^{T-1}\frac{1}{\sqrt{G_{k+1}+\epsilon}}
\le
\sqrt{T\,\frac{G_T}{\epsilon}}
= \frac{\sqrt{T\,G_T}}{\sqrt{\epsilon}} .
\]
Multiplying by $\alpha$ gives \eqref{5}.
\end{proof}

\section*{Main Proof (Step‑by‑Step Derivation)}
We prove Theorem~\ref{thm:main}.  
All expectations are taken w.r.t. the randomness of the stochastic gradients
conditioned on the past iterates.

\noindent\textbf{[Step 1]} Apply Lemma~\ref{lem:smooth} with $d=-\eta_k g_k$:
\[
f(w_{k+1})\le f(w_k)-\eta_k\inner{\nabla f(w_k)}{g_k}
+\frac{L}{2}\eta_k^{2}g_k^{2}. \tag{6}
\]

\noindent\textbf{[Step 2]} Take conditional expectation $\E_{k}[\cdot]
:=\E[\cdot\mid w_k]$ on both sides of \eqref{6}.
Using unbiasedness (Assumption~\ref{ass:unbiased}) and the definition of variance,
\[
\E_{k}[ \inner{\nabla f(w_k)}{g_k} ]
= \|\nabla f(w_k)\|^{2},
\qquad
\E_{k}[ g_k^{2} ] = \|\nabla f(w_k)\|^{2} + \E_{k}[\|g_k-\nabla f(w_k)\|^{2}]
\le \|\nabla f(w_k)\|^{2}+\sigma^{2}. \tag{7}
\]

\noindent\textbf{[Step 3]} Substitute \eqref{7} into the expectation of \eqref{6}:
\[
\E_{k}[f(w_{k+1})] \le f(w_k)
- \eta_k\|\nabla f(w_k)\|^{2}
+ \frac{L}{2}\eta_k^{2}\bigl(\|\nabla f(w_k)\|^{2}+\sigma^{2}\bigr). \tag{8}
\]

\noindent\textbf{[Step 4]} Rearrange \eqref{8} to isolate the gradient term:
\[
\eta_k\|\nabla f(w_k)\|^{2}
\le f(w_k)-\E_{k}[f(w_{k+1})]
+ \frac{L}{2}\eta_k^{2}\|\nabla f(w_k)\|^{2}
+ \frac{L}{2}\eta_k^{2}\sigma^{2}. \tag{9}
\]

\noindent\textbf{[Step 5]} Move the $\frac{L}{2}\eta_k^{2}\|\nabla f(w_k)\|^{2}$ to the left
and factor:
\[
\Bigl(\eta_k-\frac{L}{2}\eta_k^{2}\Bigr)\|\nabla f(w_k)\|^{2}
\le f(w_k)-\E_{k}[f(w_{k+1})] + \frac{L}{2}\eta_k^{2}\sigma^{2}. \tag{10}
\]

\noindent\textbf{[Step 6]} Since $0<\eta_k\le \alpha/\sqrt{\epsilon}$,
choose $\alpha$ such that $\eta_k\le 1/L$ (e.g. $\alpha\le \sqrt{\epsilon}/L$).  
Then $\eta_k-\frac{L}{2}\eta_k^{2}\ge \frac{\eta_k}{2}$, yielding
\[
\frac{\eta_k}{2}\,\|\nabla f(w_k)\|^{2}
\le f(w_k)-\E_{k}[f(w_{k+1})] + \frac{L}{2}\eta_k^{2}\sigma^{2}. \tag{11}
\]

\noindent\textbf{[Step 7]} Take full expectation over all randomness and sum
\eqref{11} from $k=0$ to $T-1$:
\[
\frac{1}{2}\sum_{k=0}^{T-1}\E[\eta_k\|\nabla f(w_k)\|^{2}]
\le f(w_0)-\E[f(w_T)] + \frac{L\sigma^{2}}{2}\sum_{k=0}^{T-1}\E[\eta_k^{2}].
\tag{12}
\]

\noindent\textbf{[Step 8]} Drop the non‑negative term $\E[f(w_T)]$ and use the
definition of $\Delta_f$:
\[
\frac{1}{2}\sum_{k=0}^{T-1}\E[\eta_k\|\nabla f(w_k)\|^{2}]
\le \Delta_f + \frac{L\sigma^{2}}{2}\sum_{k=0}^{T-1}\E[\eta_k^{2}]. \tag{13}
\]

\noindent\textbf{[Step 9]} Upper‑bound the adaptive steps:
\[
\eta_k = \frac{\alpha}{\sqrt{G_{k+1}+\epsilon}}\le \frac{\alpha}{\sqrt{\epsilon}},
\qquad
\eta_k^{2}\le \frac{\alpha^{2}}{G_{k+1}+\epsilon}
\le \frac{\alpha^{2}}{\epsilon}. \tag{14}
\]

\noindent\textbf{[Step 10]} Plug \eqref{14} into \eqref{13}:
\[
\frac{1}{2}\sum_{k=0}^{T-1}\E[\eta_k\|\nabla f(w_k)\|^{2}]
\le \Delta_f + \frac{L\sigma^{2}\alpha^{2}}{2\epsilon}\,T. \tag{15}
\]

\noindent\textbf{[Step 11]} Use Lemma~\ref{lem:stepsum} (inequality \eqref{5}) to lower‑bound the left‑hand side:
\[
\sum_{k=0}^{T-1}\E[\eta_k\|\nabla f(w_k)\|^{2}]
\ge
\Bigl(\min_{0\le k<T}\E[\|\nabla f(w_k)\|^{2}]\Bigr)
\sum_{k=0}^{T-1}\E[\eta_k].
\]
Taking expectations inside the sum and applying \eqref{5} gives
\[
\sum_{k=0}^{T-1}\E[\eta_k]\;\le\;
\frac{2\alpha}{\sqrt{\epsilon}}\sqrt{T\,\E[G_T] } .
\tag{16}
\]

\noindent\textbf{[Step 12]} Since $G_T=\sum_{i=0}^{T-1}g_i^{2}\le T G_{\max}^{2}$,
\[
\sqrt{T\,\E[G_T]} \le \sqrt{T\cdot T G_{\max}^{2}} = T G_{\max}.
\tag{17}
\]

\noindent\textbf{[Step 13]} Combine \eqref{15}, \eqref{16}, and \eqref{17}:
\[
\frac{1}{2}\,
\Bigl(\min_{k<T}\E[\|\nabla f(w_k)\|^{2}]\Bigr)
\frac{2\alpha}{\sqrt{\epsilon}}\,T G_{\max}
\;\le\;
\Delta_f + \frac{L\sigma^{2}\alpha^{2}}{2\epsilon}\,T .
\]
Cancel the factor $2$ and divide by $T$:
\[
\min_{k<T}\E[\|\nabla f(w_k)\|^{2}]
\;\le\;
\frac{\sqrt{\epsilon}\,\Delta_f}{\alpha\,G_{\max}\,T}
\;+\;
\frac{L\sigma^{2}\alpha}{\sqrt{\epsilon}\,G_{\max}} .
\tag{18}
\]

\noindent\textbf{[Step 14]} For a cleaner bound we replace $G_{\max}$ by
$G_{\max}\le \sqrt{\epsilon}+G_{\max}$ and absorb constants, arriving at the
simplified expression (4) stated in Theorem~\ref{thm:main}:
\[
\frac{1}{T}\sum_{k=0}^{T-1}\E[\|\nabla f(w_k)\|^{2}]
\le
\frac{2\Delta_f}{\alpha\sqrt{T\epsilon}}
+
\frac{2L G_{\max}^{2}\alpha}{\sqrt{T\epsilon}} .
\]
$\square$

\section*{Rate Derivation (With Constants)}
From \eqref{4},
\[
\mathcal{O}\!\left(\frac{1}{\sqrt{T}}\right)
\quad\text{with constant }
C = \frac{2\Delta_f}{\alpha\sqrt{\epsilon}} + 2L G_{\max}^{2}\alpha .
\]
Choosing $\alpha = \sqrt{\Delta_f/(L G_{\max}^{2}\epsilon)}$ balances the two terms,
yielding the optimal constant
\[
C_{\text{opt}} = 4\sqrt{L\,\Delta_f}\,G_{\max}/\sqrt{\epsilon}.
\]

\section*{Special Cases}
\begin{itemize}
\item \textbf{Deterministic gradients ($\sigma=0$).}
  The second term in \eqref{4} disappears and the bound reduces to
  $2\Delta_f/(\alpha\sqrt{T\epsilon})$, i.e. a pure $1/\sqrt{T}$ decay.
\item \textbf{Large initial $\epsilon$.}
  If $\epsilon$ dominates $G_T$, then $\eta_k\approx \alpha/\sqrt{\epsilon}$ is
  essentially constant and the method behaves like standard SGD with stepsize
  $\alpha/\sqrt{\epsilon}$, recovering the classic $O(1/\sqrt{T})$ SGD rate.
\item \textbf{Coordinate‑wise extension.}
  For $d$‑dimensional $w$, replace $G_{k+1}$ by the vector of
  accumulated squared coordinates; the same proof holds component‑wise,
  leading to $\mathcal{O}(\frac{d}{\sqrt{T}})$ in the worst case.
\end{itemize}

\end{document}